\scanpage{83}
% archive foliation: 82
% lines 1--3
\deleted{$X_0\unsure{/S}$} d'une situation $X_0/I_0/\deleted{k_0}k_0$
où $k_0$ est un sous-corps. [\unsure{OPS} \unsure{lors} que
$k$ \unsure{soit|est} une clôture alg. de $k_0$\ldots].
% lines 4--8
Alors $\Gamma=\operatorname{Gal}(\overleftrightarrow{k/k_0})$ opère sur
$\Sigma^{ext}$ par transport de structure, donc on a
une \deleted{la} opération aussi sur $P(X/I/k)$,
\deleted{donc} \inserted{ou encore} un hom
\[
\Gamma\longrightarrow {}^{P(X/I/k)}\widetilde\Gamma_\nu\simeq\operatorname{Aut}(\Sigma^{ext})
\]
% lines 9--11
qui par image inverse \inserted{d'extensions} donne naissance
à l'extension
\diagram{\begin{tikzcd}[ampersand replacement=\&,column sep=small]
1\arrow[r,dash]\&\pi_1(X\setminus I,\unsure{a|x})\arrow[r,dash]\&\pi_1(X_0\setminus I_0,\unsure{a_0|x_0})\arrow[r]\&\pi_1(k_0\unsure{/|,}k)\arrow[r,dash]\&1
\end{tikzcd}}
% lines 12--13
\ldots\quad $\Gamma\xrightarrow{\chi}\operatorname{Aut}(T)\simeq\widehat{\mathbb Z}^{*}$ est l'hom de
Kummer-Tate\ldots
% lines 14--16
\textcircled{II} Plus précisément supposons $k_0=Q$, $X_0=\mathbb P^1_Q$,
\inserted{$\overline Q$ clôture alg de $Q$ dans $C$,}
\[
I=\mu_{\nu+1}(\overline Q)=\mu_{\nu+1}(C^*)\subset\overline Q=\unsure{\mathbb E^1_Q|\mathbb A^1_Q}(\overline Q)\subset\mathbb P^1_Q(\overline Q
\]
% lines 17--21
On trouve un isomorphisme canonique,
via l'exponentielle et les lacets
\[
\deleted{I\simeq[0,\unsure{\nu|2}]},
\qquad \Sigma(X/I/k)^{ext}\simeq L_\nu^{ext}
\]
donc\quad $P(X,I/k)=$ torseur trivial
